%%%%%%%%%%%%%%%%%%%%%%%%%%%%%%%%%%%%%%%%%%%%%%%%%%%%%%%%%
%%%%% Nouvelle commande \domino[Sup][Inf](Coord sw) %%%%%
%%%%%%%%%%%%%%%%%%%%%%%%%%%%%%%%%%%%%%%%%%%%%%%%%%%%%%%%%

\definecolor{color_domino}{HTML}{FFE9C0}
\newcommand\domino{}
\def\domino[#1][#2](#3){
	
	%%% Aspect du domino %%%

	\draw[rounded corners , line width = 1pt , fill = color_domino] ($(#3)+(0,0)$) rectangle ($(#3)+(2,3)$) ;
	\draw[fill , line width = 1pt , color = black!70!white] ($(#3)+(0.25,1.5)$) -- ($(#3)+(1,1.5)$) circle[radius=0.075] -- ($(#3)+(1.75,1.5)$) ;
	
	%%% Compteur de points partie supérieure %%%
	
	\ifnum #1 = 1
		\draw[fill] ($(#3)+(1,2.25)$) circle[radius=0.125] ;
	\fi
	
	\ifnum #1 = 2
		\draw[fill] ($(#3)+(0.5,2.625)$) circle[radius=0.125] ;
		\draw[fill] ($(#3)+(1.5,1.875)$) circle[radius=0.125] ;
	\fi
	
	\ifnum #1 = 3
		\draw[fill] ($(#3)+(0.5,2.625)$) circle[radius=0.125] ;
		\draw[fill] ($(#3)+(1.5,1.875)$) circle[radius=0.125] ;
		\draw[fill] ($(#3)+(1,2.25)$) circle[radius=0.125] ;
	\fi
	
	\ifnum #1 = 4
		\draw[fill] ($(#3)+(0.5,2.625)$) circle[radius=0.125] ;
		\draw[fill] ($(#3)+(1.5,1.875)$) circle[radius=0.125] ;
		\draw[fill] ($(#3)+(1.5,2.625)$) circle[radius=0.125] ;
		\draw[fill] ($(#3)+(0.5,1.875)$) circle[radius=0.125] ;
	\fi
	
	\ifnum #1 = 5
		\draw[fill] ($(#3)+(0.5,2.625)$) circle[radius=0.125] ;
		\draw[fill] ($(#3)+(1.5,1.875)$) circle[radius=0.125] ;
		\draw[fill] ($(#3)+(1.5,2.625)$) circle[radius=0.125] ;
		\draw[fill] ($(#3)+(0.5,1.875)$) circle[radius=0.125] ;
		\draw[fill] ($(#3)+(1,2.25)$) circle[radius=0.125] ;
	\fi
	
	\ifnum #1 = 6
		\draw[fill] ($(#3)+(0.5,2.625)$) circle[radius=0.125] ;
		\draw[fill] ($(#3)+(1.5,1.875)$) circle[radius=0.125] ;
		\draw[fill] ($(#3)+(1.5,2.625)$) circle[radius=0.125] ;
		\draw[fill] ($(#3)+(0.5,1.875)$) circle[radius=0.125] ;
		\draw[fill] ($(#3)+(1.5,2.25)$) circle[radius=0.125] ;
		\draw[fill] ($(#3)+(0.5,2.25)$) circle[radius=0.125] ;
	\fi
	
	
	%%% Compteur de points partie inférieure %%%
	
	\ifnum #2 = 1
		\draw[fill] ($(#3)+(1,0.75)$) circle[radius=0.125] ;
	\fi
	
	\ifnum #2 = 2
		\draw[fill] ($(#3)+(0.5,1.125)$) circle[radius=0.125] ;
		\draw[fill] ($(#3)+(1.5,0.375)$) circle[radius=0.125] ;
	\fi
	
	\ifnum #2 = 3
		\draw[fill] ($(#3)+(0.5,1.125)$) circle[radius=0.125] ;
		\draw[fill] ($(#3)+(1.5,0.375)$) circle[radius=0.125] ;
		\draw[fill] ($(#3)+(1,0.75)$) circle[radius=0.125] ;
	\fi
	
	\ifnum #2 = 4
		\draw[fill] ($(#3)+(0.5,1.125)$) circle[radius=0.125] ;
		\draw[fill] ($(#3)+(1.5,0.375)$) circle[radius=0.125] ;
		\draw[fill] ($(#3)+(1.5,1.125)$) circle[radius=0.125] ;
		\draw[fill] ($(#3)+(0.5,0.375)$) circle[radius=0.125] ;
	\fi
	
	\ifnum #2 = 5
		\draw[fill] ($(#3)+(0.5,1.125)$) circle[radius=0.125] ;
		\draw[fill] ($(#3)+(1.5,0.375)$) circle[radius=0.125] ;
		\draw[fill] ($(#3)+(1.5,1.125)$) circle[radius=0.125] ;
		\draw[fill] ($(#3)+(0.5,0.375)$) circle[radius=0.125] ;
		\draw[fill] ($(#3)+(1,0.75)$) circle[radius=0.125] ;
	\fi
	
	\ifnum #2 = 6
		\draw[fill] ($(#3)+(0.5,1.125)$) circle[radius=0.125] ;
		\draw[fill] ($(#3)+(1.5,0.375)$) circle[radius=0.125] ;
		\draw[fill] ($(#3)+(1.5,1.125)$) circle[radius=0.125] ;
		\draw[fill] ($(#3)+(0.5,0.375)$) circle[radius=0.125] ;
		\draw[fill] ($(#3)+(1.5,0.75)$) circle[radius=0.125] ;
		\draw[fill] ($(#3)+(0.5,0.75)$) circle[radius=0.125] ;
	\fi
}




%%%%%%%%%%%%%%%%%%%%%%%%%%%%%%%%%%%%%%%%%%%%%%
%%%%% Macros pour les diagrammes de Venn %%%%%
%%%%%%%%%%%%%%%%%%%%%%%%%%%%%%%%%%%%%%%%%%%%%%

% Rectangle de bordure
\def\Rect[#1]{
	\ifnum #1 = 2
		(0,0) rectangle (6,4)
	\fi
	\ifnum #1 = 3
		(0,0) rectangle (6,5.732)
	\fi	
}


%%% Pour un diagramme de Venn à deux cercles %%%

% Cercle de A
\def\A{(2,2)  circle (1.4142)}

% Cercle de B
\def\B{(4,2) circle[radius=1.4142]}

% A union B
\def\AouB{(3,1) arc(-135:135:1.4142) arc (45:315:1.4142)}

% A inter B
\def\AetB{(3,1) arc(-45:45:1.4142) arc (135:225:1.4142)}

% A privé de B
\def\AmB{(3,1) arc(315:45:1.4142) arc (135:225:1.4142)}

% B privé de A
\def\BmA{(3,1) arc(-135:135:1.4142) arc (45:-45:1.4142)}


%%% Pour un diagramme de Venn à trois cercles %%%

% Cercle de C
\def\C{(3,3.732) circle[radius=1.4142]}

% A union B union C
\def\AouBouC{(3,1) arc(-135:75:1.4142) arc (-15:195:1.4142) arc (105:315:1.4142)}

% A inter B inter C
\def\AetBetC{(3,3) arc(135:165:1.4142) arc (255:285:1.4142) arc (15:45:1.4142)}

% A union C
\def\AouC{(2,2)+(15:1.4142) arc(-75:195:1.4142) arc (105:375:1.4142)}

% A inter C
\def\AetC{(2,2)+(15:1.4142) arc(-75:-165:1.4142) arc (105:15:1.4142)}

% A privé de C
\def\AmC{(2,2)+(15:1.4142) arc(-75:-165:1.4142) arc (105:375:1.4142)}

% C privé de A
\def\CmA{(2,2)+(15:1.4142) arc(-75:195:1.4142) arc (105:15:1.4142)}

% B union C
\def\BouC{(4,2)+(165:1.4142) arc(-195:75:1.4142) arc (-15:255:1.4142)}

% B inter C
\def\BetC{(4,2)+(165:1.4142) arc(165:75:1.4142) arc (-15:-105:1.4142)}

% B privé de C
\def\BmC{(4,2)+(165:1.4142) arc(-195:75:1.4142) arc (-15:-105:1.4142)}

% C privé de B
\def\CmB{(4,2)+(165:1.4142) arc(165:75:1.4142) arc (-15:255:1.4142)}

% A seul
\def\Aseul{(3,1) arc(-45:-255:1.4142) arc (-165:-105:1.4142) arc (165:225:1.4142)}

% B seul
\def\Bseul{(3,1) arc(-135:75:1.4142) arc (-15:-75:1.4142) arc (15:-45:1.4142)}

% C seul
\def\Cseul{(4,2)+(75:1.4142) arc(-15:195:1.4142) arc (105:45:1.4142) arc (135:75:1.4142)}

% A inter B moins C
\def\AetBmC{(3,1) arc(-45:15:1.4142) arc (-75:-105:1.4142) arc (165:225:1.4142)}

% A inter C moins B
\def\AetCmB{(3,3) arc(45:105:1.4142) arc (-165:-105:1.4142) arc (165:135:1.4142)}

% B inter C moins A
\def\BetCmA{(3,3) arc(135:75:1.4142) arc (-15:-75:1.4142) arc (15:45:1.4142)}




